\chapter*{附\ \ 录}
\addcontentsline{toc}{chapter}{附\ \ 录}

\setcounter{section}{0}
\renewcommand{\thesection}{\Alph{section}}
\renewcommand{\theHsection}{appendix.\Alph{section}}
\renewcommand{\theHsubsection}{\theHsection.\arabic{subsection}}
\counterwithin{figure}{section}
\counterwithin{table}{section}
\counterwithin{equation}{section}
\renewcommand{\thefigure}{\thesection-\arabic{figure}}
\renewcommand{\thetable}{\thesection-\arabic{table}}
\renewcommand{\theequation}{\thesection-\arabic{equation}}
\renewcommand{\theHfigure}{\theHsection.\arabic{figure}}
\renewcommand{\theHtable}{\theHsection.\arabic{table}}
\renewcommand{\theHequation}{\theHsection.\arabic{equation}}

\section{仿真环境与验证方法}

\subsection{工具配置}

{[}待填写：仿真器、RISC-V编译工具、参考模型及依赖工具的版本，源码提交、配置参数、运行命令和日志保存位置。{]}

\subsection{测试平台}

{[}待填写：单元、模块和系统级测试的划分，测试台结构、时钟复位、输入生成方式及输出采集方式；说明仿真覆盖范围，不将仿真结果视为板级电气或实物验证结论。{]}

\subsection{判定标准}

{[}待填写：各类测试的验证目的、参考模型、期望行为、数值误差容限、超时条件和通过标准，以及波形、数值结果和验证结论的记录方式。{]}

\section{CPU仿真验证}

\subsection{指令集验证}

{[}待填写：支持的RV32指令集能力与扩展配置、官方架构测试及定向用例，未启用扩展的非法指令行为、测试结果和已知限制。{]}

\subsection{流水线控制}

{[}待填写：译码、发射和提交行为，数据冒险、旁路、暂停和冲刷的验证用例，关键流水线信号、退休结果及判定依据。{]}

\subsection{指令提取}

{[}待填写：取指时序、压缩指令处理、分支预测和重定向的测试输入与关键波形；验证压缩指令配置下预测查询、更新及相关性能计数器的行为。{]}

\subsection{指令执行}

{[}待填写：整数运算、移位、比较和乘除法的测试配置，符号组合、边界输入及除法特殊情况，与参考结果的逐项比较和执行完成时序。{]}

\subsection{中断和异常}

{[}待填写：已实现CSR的读写行为、异常触发、中断响应和返回流程，异常原因与返回地址、流水线处理及关键波形的判定结果。{]}

\subsection{存储访问}

{[}待填写：缓存命中与未命中、读写访问、字节使能、访问对齐和PMA边界的验证用例，访存时序、读回数据及异常行为。{]}

\section{NPU仿真验证}

\subsection{计算正确性}

{[}待填写：PE和脉动阵列的测试输入、INT8乘法与INT32累加约定，符号组合、累加边界、参考结果及逐项数值比较。{]}

\subsection{分块计算}

{[}待填写：矩阵规模、数据布局、分块参数和尾块处理的测试组合，输入输出地址、块间拼接及与完整矩阵参考计算的比较结果。{]}

\subsection{DMA传输}

{[}待填写：DMA读写搬运、传输长度、地址边界、突发传输和背压场景，握手波形、内存读回数据及传输完成结果。{]}

\subsection{任务控制}

{[}待填写：寄存器配置、任务启动、忙闲状态、完成通知和连续任务提交的测试序列，以及异常输入、背压和任务重新启动时的状态转换。{]}

\subsection{可选功能}

{[}待填写：卷积展开、定点后处理等已实现可选功能的测试输入、参考结果和边界情况；分别记录开启与关闭配置下的功能行为、合法响应及回退路径。{]}

\section{系统互连与存储仿真验证}

\subsection{总线协议}

{[}待填写：总线各通道握手、突发传输、字节使能和响应行为的测试配置，背压场景、关键协议波形及通过标准。{]}

\subsection{访问仲裁}

{[}待填写：多主设备并发访问、持续请求和不同背压条件下的测试序列，仲裁顺序、请求完成情况及背压传播结果。{]}

\subsection{地址译码}

{[}待填写：地址映射、存储与外设访问边界、非法地址及禁用模块访问的测试用例，目标选择、读写结果和错误响应。{]}

\subsection{共享数据访问}

{[}待填写：CPU、NPU DMA和图像链路的共享数据访问场景，缓冲区交接、缓存维护、数据可见性及读写结果的验证；按实际访问约定检查，不假定硬件自动维持缓存一致性。{]}

\section{图像采集与显示链路仿真验证}

\subsection{图像采集}

{[}待填写：图像输入时序、帧同步、行同步、像素有效信号、像素格式和帧标识的测试配置，采集数据及帧完整性的判定结果。{]}

\subsection{输入预处理}

{[}待填写：已实现的像素格式转换、尺寸处理和模型输入组织，测试图像、输入输出尺寸、通道排列及与参考处理结果的比较。{]}

\subsection{跨时钟域}

{[}待填写：不同采集、处理和显示时钟条件下的跨域同步，数据与控制信号传递、复位恢复及边界时序的测试波形和判定结果。{]}

\subsection{帧缓存}

{[}待填写：帧缓存读写、缓冲区切换、所有权交接和读写冲突的测试场景，帧完整性、数据保持及冲突处理结果。{]}

\subsection{显示输出}

{[}待填写：显示扫描、像素输出、同帧结果配对和输入输出一致性的仿真配置，输出时序、参考像素及显示数据核对结果。{]}

\section{系统联合仿真验证}

\subsection{完整模型推理}

{[}待填写：模型版本、输入尺寸、算子执行后端及测试输入，CPU与NPU协同、端到端数据流和推理完成情况，运行日志及验证边界。{]}

\subsection{连续帧处理}

{[}待填写：多帧输入顺序、帧标识、任务衔接和缓冲区交接的测试序列，连续运行中的完成通知、帧配对及结果保持情况。{]}

\subsection{运行配置}

{[}待填写：加速启用与禁用、CPU回退及显示配置组合的测试矩阵，各配置的运行路径、任务完成情况和结果一致性；区别于单个NPU可选功能的验证。{]}

\subsection{结果核对}

{[}待填写：整网输出与软件参考结果的比较、数值误差统计、帧标识及结果配对检查，参考图像、核对日志和通过结论；注明误差容限与未覆盖项目。{]}
